\section{Experiments}
\label{sec:experiments}

\subsection{Evaluation Benchmark and Protocol}
\label{sec:evaluation_protocol}

We evaluate personalized designer role-play in a cross-scenario transfer
setting using behavioral data from 162 designers. Information collected in
Questions 1--6 (Q1--Q6) is used to construct each participant's professional
background, 30-dimensional value profile (V30), and historical response to a
design problem. Question 7 (Q7) presents a previously unseen wicked-design
dilemma involving conflicts among commercial efficiency, accessibility, and
manipulative recommendation. The target output is the participant's response
to Q7.

Q7 contains three elicitation stages: an intuitive first reaction
(\textbf{Seg1}), a decision-weighing reasoning process (\textbf{Seg2}), and a
final design bottom line (\textbf{Seg3}). Our main experiment focuses on Seg2,
where both a structured C--K trajectory and its natural-language realization
are available for evaluation. Seg1 and Seg3 are used only in auxiliary
analyses. All main tables report the full cohort of 162 participants.

The first ten participants were used during prompt-format inspection and
boundary-error diagnosis. We therefore distinguish configuration selection
from cross-backbone evaluation: GPT-4o is used as the development backbone for
selecting the final prompt configuration, after which the configuration is
frozen and transferred unchanged to the remaining backbones. We report this
protocol explicitly because the full-cohort results should be interpreted as
model-transfer evidence rather than as a strictly untouched participant-level
test set.

\subsection{Compared Methods}
\label{sec:baselines}

We compare our method with three prompt-only baselines and two trainable role-playing
baselines. Prompt-only baselines are evaluated on Qwen3-8B,
Llama-3.1-8B-Instruct, Gemini 3 Flash, and GPT-4o. CoSER and RAR-RIA are
adapted using Qwen3-8B and Llama-3.1-8B-Instruct. The focused main table
reports Qwen3-8B, Llama-3.1-8B-Instruct, and GPT-4o; the remaining
cross-backbone results are reported in the supplementary material.

\paragraph{Prompt-only baselines.}
All prompt-only baselines generate both a C--K state sequence and a
natural-language answer in one response. They share the same Q7 situation and
decoding configuration.

\textbf{SP} uses only the Q7 situation and a generic designer role, without
participant-specific information.

\textbf{PP} additionally provides the target participant's personal and
professional information (PI), including design field, educational
background, work context, project role, and design experience.

\textbf{ICL-Q2} extends PP with three Q2 demonstrations from other
participants. Each demonstration contains PI, a Q2 C--K state sequence, and a
Q2 natural-language response. The target participant is always excluded from
the demonstration set.

\paragraph{Trainable role-playing baselines.}
We adapt CoSER~\citep{wang2025coser} and RAR-RIA~\citep{tang2026thinking} to
the designer questionnaire data and train them on Q1--Q6 using their
respective LoRA supervised fine-tuning procedures. CoSER uses its
\texttt{model\_speech} field as the prediction, whereas RAR-RIA uses its
\texttt{model\_answer} field. These methods do not define an explicit C--K
trajectory; consequently, only natural-language and value-expression metrics
are applicable to them.

\subsection{Configuration of the Proposed Method}
\label{sec:proposed_configuration}

Our method performs trait-conditioned aligned C--K generation in a single
model call. Given the public Q7
situation, it first identifies four situation-relevant dimensions
(\emph{active4}) from the target participant's V30 profile. The prompt contains
(i) three other-participant Q7 Seg2 examples used to calibrate C--K operation
boundaries and output granularity, (ii) target PI, (iii) the four activated
V30 dimensions with their target-specific values and interpretations,
(iv) the target participant's complete Q2 state sequence and response, and
(v) selected trait-conditioned rules learned from Q1--Q6, and (vi) the full
Q7 situation and Seg2 instruction. The target participant's Q7 answer is
never included in the prompt.

The model returns a single JSON object:

\begin{equation}
\mathcal{U}_i =
\left[(m_{i1},y_{i1}),\ldots,(m_{iT_i},y_{iT_i})\right],
\end{equation}

where $m_{it}$ is one of
$\{\mathrm{K},\mathrm{C},\mathrm{K{\to}C},\mathrm{C{\to}K},
\mathrm{K{\to}K},\mathrm{C{\to}C}\}$ and $y_{it}$ is the text realizing that
operation. The trajectory and final response are deterministic projections:

\begin{equation}
\hat{M}_i=[m_{i1},\ldots,m_{iT_i}], \qquad
\hat{Y}_i=y_{i1}\oplus\cdots\oplus y_{iT_i}.
\end{equation}

Thus, our method does not use a separately called planner or realizer. It also does
not generate one state per model call. All aligned units are produced in the
same context and the same model call.

The proposed configuration does not place the complete
30-dimensional profile in the prompt. The full profile and active4 provide
overlapping signals and jointly reduced both trajectory discrimination and
value-expression alignment in the development experiment. It therefore uses
active4 as a compact, situation-conditioned projection of V30.

The proposed method is evaluated on Qwen3-8B, Llama-3.1-8B-Instruct, GPT-4o,
and Claude Sonnet 4.6. GPT-4o is the development and ablation backbone; the
other three models evaluate transfer of the frozen
configuration. We use temperature $0.7$, top-$p$ $0.9$, and at most 1,024 new
tokens for generation.

\subsection{Evaluation Metrics}
\label{sec:evaluation_metrics}

We evaluate four complementary properties: response alignment, trajectory
fidelity, participant-specific personalization, and resistance to population-
level collapse. Metrics are calculated on Seg2.

\paragraph{Response alignment.}
\textbf{BGE own-user margin} compares the BGE-M3 similarity of a generated
response to its target participant's gold response against its average
similarity to other participants' gold responses. A positive value indicates
preferential alignment with the intended participant.

\paragraph{Trajectory fidelity.}
\textbf{KC-NES} is one minus the normalized edit distance between the
generated and gold C--K sequences. We additionally report exact match
(\textbf{EM}) and absolute sequence-length error (\textbf{LenErr}) as
diagnostic measures.

\paragraph{Trajectory personalization.}
\textbf{KC own-user AUC} measures how often the predicted trajectory is more
similar to its own participant's gold trajectory than to another
participant's trajectory. \textbf{KC own-user margin} measures the
corresponding mean similarity difference. These metrics distinguish genuine
participant conditioning from a population-level trajectory that is
moderately similar to everyone.

\paragraph{Value-expression faithfulness.}
A frozen evaluator extracts directional V30 claims from human and generated
responses. \textbf{Human-profile Expression Precision, Recall, and F1}
compare generated claims with claims expressed in the same participant's
human Q7 response. \textbf{Profile-polarity precision} measures whether
generated claim directions agree with the participant's V30 profile.
\textbf{Expression coverage} reports the fraction of V30 dimensions expressed
in the response and is treated as descriptive rather than intrinsically
better when larger. \textbf{Own-profile advantage} is the generated
response's alignment with its intended V30 profile minus its average
alignment with other participants' profiles. The same frozen extraction
procedure is applied to human and model responses.

\paragraph{Anti-collapse.}
\textbf{KCDR} measures how much of the pairwise diversity in human C--K
trajectories is retained by generated trajectories. \textbf{TGC} measures the
agreement between pairwise participant geometry in generated and gold
trajectory spaces. We also report the number of unique predicted trajectories
(\textbf{Unique}) and the frequency of the most common exact trajectory
(\textbf{Top share}). Unique and Top share are interpreted jointly with
KC-NES, KCDR, and TGC because unconstrained randomness can increase diversity
without improving fidelity.

\begin{table}[t]
\centering
\small
\caption{Evaluation dimensions and primary metrics. Diagnostic metrics are
reported in the appendix.}
\label{tab:evaluation_metrics}
\begin{tabular}{p{0.22\linewidth}p{0.68\linewidth}}
\toprule
\textbf{Dimension} & \textbf{Metrics} \\
\midrule
Response alignment & BGE own-user margin $\uparrow$ \\
Trajectory fidelity & KC-NES $\uparrow$ \\
Trajectory personalization & KC own-user AUC $\uparrow$; KC own-user margin $\uparrow$ \\
Value faithfulness & Human-profile Expression F1 $\uparrow$; profile-polarity precision $\uparrow$; own-profile advantage $\uparrow$ \\
Anti-collapse & KCDR $\uparrow$; TGC $\uparrow$; Top share $\downarrow$ \\
\bottomrule
\end{tabular}
\end{table}

\subsection{Main Results}
\label{sec:main_results}

\subsubsection{Comparison with Prompt and Trainable Baselines}

\begin{table*}[t]
\centering
\small
\caption{Personalization and anti-collapse results on 162 participants.
Higher is better for BGE, KC-NES, Polarity, Own-profile advantage, KCDR, and
Unique; lower is better for LenErr and TopShare. A value is bolded only when
Ours is the best-performing method within that backbone; SP serves only as a
non-personalized reference and is excluded from this comparison. ``--''
denotes an unavailable or non-applicable metric.}
\label{tab:r5_main_results}
\resizebox{\textwidth}{!}{
\begin{tabular}{llrrrrrrrr}
\toprule
\textbf{Model} & \textbf{Method} &
\textbf{BGE} &
\textbf{KC-NES} &
\textbf{LenErr} &
\textbf{Polarity} &
\textbf{Own-profile Adv.} &
\textbf{KCDR} &
\textbf{Unique} &
\textbf{Top share} \\
\midrule
Qwen3-8B & SP       & $-0.0000$ & 0.2977 & 1.4444 & -- & -- & 0.0000 & 1 & 1.0000 \\
         & PP       & $-0.0000$ & 0.2773 & 2.0370 & -- & -- & 0.3953 & 20 & 0.3889 \\
         & ICL-Q2   & $-0.0017$ & 0.2435 & 3.6173 & -- & -- & 0.3811 & 47 & 0.2037 \\
         & CoSER    & $-0.0056$ & -- & -- & 0.4940 & $-0.0296$ & -- & -- & -- \\
         & RAR-RIA  & $-0.0075$ & -- & -- & 0.5141 & 0.0157 & -- & -- & -- \\
         & \textbf{Ours} & \textbf{0.0035} & 0.2521 & \textbf{1.6049} & \textbf{0.5264} & \textbf{0.0432} & \textbf{0.5472} & \textbf{50} & \textbf{0.1173} \\
\midrule
Llama-3.1-8B & SP       & $-0.0000$ & 0.2522 & 4.3210 & -- & -- & 0.0000 & 1 & 1.0000 \\
             & PP       & 0.0004 & 0.2550 & 7.6852 & -- & -- & 0.7783 & 154 & 0.0123 \\
             & ICL-Q2   & $-0.0073$ & 0.2589 & 5.0185 & -- & -- & 0.7580 & 156 & 0.0123 \\
             & CoSER    & $-0.0035$ & -- & -- & 0.5370 & 0.0491 & -- & -- & -- \\
             & RAR-RIA  & 0.0040 & -- & -- & 0.4406 & $-0.0309$ & -- & -- & -- \\
             & \textbf{Ours} & $-0.0013$ & 0.2202 & \textbf{2.4198} & 0.5353 & 0.0485 & 0.7002 & 119 & 0.1728 \\
\midrule
GPT-4o & SP       & 0.0030 & 0.2764 & 1.5926 & -- & -- & 0.3482 & 40 & 0.4506 \\
       & PP       & $-0.0012$ & 0.2650 & 1.5679 & -- & -- & 0.3908 & 43 & 0.3457 \\
       & ICL-Q2   & 0.0003 & 0.2504 & 2.3642 & -- & -- & 0.5888 & 121 & 0.0494 \\
       & \textbf{Ours} & $-0.0017$ & 0.2638 & 3.8272 & 0.5298 & 0.0210 & \textbf{0.5904} & \textbf{132} & \textbf{0.0309} \\
\bottomrule
\end{tabular}}
\end{table*}

On Qwen3-8B, our method achieves the strongest V30 alignment, with the highest
profile-polarity precision (\(0.5264\)) and own-profile advantage
(\(0.0432\)). It also produces 50 unique trajectories and reduces TopShare
from \(1.0000\) under SP to \(0.1173\). On GPT-4o, it produces 132 unique
trajectories, lowers TopShare to \(0.0309\), and achieves the highest KCDR
(\(0.5904\)). These results show that the proposed conditioning scheme
reduces population-level trajectory collapse while retaining alignment with
the intended participant's value profile.

\subsection{Ablation Analysis}
\label{sec:ablation}

Our method uses PI, the target participant's Q2 record, Q7-specific active4
cues, selected trait-conditioned rules, and other-participant aligned Q7
examples. Table~\ref{tab:r5_ablation_design} defines ablations relative to
this final configuration.

\begin{table}[t]
\centering
\small
\caption{Component ablations of the proposed method. All variants retain the same Q7 task,
aligned-unit output contract, Q7 examples, and decoding parameters.}
\label{tab:r5_ablation_design}
\begin{tabular}{lccccc}
\toprule
\textbf{Variant} & \textbf{PI} & \textbf{Q2} & \textbf{active4} &
\textbf{Rules} & \textbf{Full30} \\
\midrule
Ours & \checkmark & \checkmark & \checkmark & \checkmark & $\times$ \\
Ours $+$ Full30 & \checkmark & \checkmark & \checkmark & \checkmark & \checkmark \\
Ours w/o Q2 & \checkmark & $\times$ & \checkmark & \checkmark & $\times$ \\
Ours w/o active4 & \checkmark & \checkmark & $\times$ & \checkmark & $\times$ \\
Ours w/o rules & \checkmark & \checkmark & \checkmark & $\times$ & $\times$ \\
Ours w/ shuffled rules & \checkmark & \checkmark & \checkmark & shuffled & $\times$ \\
\bottomrule
\end{tabular}
\end{table}

\begin{table*}[t]
\centering
\small
\caption{Ablation results on GPT-4o with 162 participants. Higher is better
for BGE, KC-NES, Polarity, Own-profile advantage, KCDR, and Unique; lower is
better for LenErr and TopShare.}
\label{tab:r5_ablation_results}
\resizebox{\textwidth}{!}{
\begin{tabular}{lrrrrrrrr}
\toprule
\textbf{Configuration} & \textbf{BGE} & \textbf{KC-NES} &
\textbf{LenErr} & \textbf{Polarity} & \textbf{Own-profile Adv.} &
\textbf{KCDR} & \textbf{Unique} & \textbf{TopShare} \\
\midrule
Ours                   & $-0.0017$ & 0.2638 & 3.8272 & 0.5298 & 0.0210 & 0.5904 & 132 & 0.0309 \\
Ours w/o rules         & $-0.0008$ & 0.2674 & 3.8580 & 0.5509 & 0.0446 & 0.5736 & 139 & 0.0247 \\
Ours \(+\) Full30      & $-0.0007$ & 0.2632 & 3.8148 & 0.5330 & 0.0241 & 0.5852 & 138 & 0.0247 \\
Ours w/o Q2            & $-0.0009$ & 0.2706 & 3.1543 & 0.5550 & 0.0542 & 0.5732 & 125 & 0.0617 \\
Ours w/o active4       & 0.0017 & 0.2688 & 3.8765 & 0.4881 & $-0.0163$ & 0.5646 & 130 & 0.0494 \\
Ours w/ shuffled rules & 0.0018 & 0.2674 & 3.6049 & 0.5252 & 0.0220 & 0.5974 & 134 & 0.0370 \\
\bottomrule
\end{tabular}}
\end{table*}

Table~\ref{tab:r5_ablation_results} shows that removing Q2 increases TopShare
from \(0.0309\) to \(0.0617\) and reduces Unique from 132 to 125, indicating
that participant history helps differentiate trajectories. Removing the
learned rules reduces KCDR from \(0.5904\) to \(0.5736\), whereas adding
Full30 provides no KCDR gain over active4-only (\(0.5852\) vs.\
\(0.5904\)). We therefore retain Q2, the rule library, and the more compact
active4-only profile representation in the final configuration.
